% GENERATED FILE. Edit canonical lessons, then run npm run generate:books.
% canonical-source-hash: fnv1a64:732b6ddc5291aad2
% canonical-lessons: FR-C252-impot, FR-C252-formulaire, FR-C252-enveloppe, FR-C252-timbre, FR-C252-colis

\chapter{Tax, Form, Envelope, Stamp, Parcel}
\label{ch:fr-tax-form-envelope-stamp-parcel}
\hlchaptermodality{252}

\begin{chapteropening}
\textbf{By the end of this chapter:} \emph{I can say a tax, a form, an envelope, a stamp and a parcel.}

Complete the last lesson of chapter 252: I can say a tax, a form, an envelope, a stamp and a parcel.
\end{chapteropening}

\section[l'impôt]{l'impôt — a tax}

\label{lesson:FR-C252-impot}



\begin{quote}
Before the new one: say the French for the truth, then the French for a lie.
\end{quote}

\subsection*{You'll want to know: l'impôt}
\textbf{l'impôt} — \textquotedblleft{}a tax\textquotedblright{}.

\begin{cousinweb}
\textbf{impôt} $\leftarrow$ Middle French \emph{impost} $\leftarrow$ Medieval Latin \emph{impositum}, \textquotedblleft{}a thing laid on\textquotedblright{}, from \emph{imponere} (EtymDB 2.1) — a learned borrowing; the circumflex marks the \emph{s} French later stopped writing.
\end{cousinweb}

\begin{grammarlens}[title={What you've built}]
One more word for this chapter.
\end{grammarlens}

\subsection*{Your turn}
\begin{itemize}
\raggedright
  \item \emph{Say it:} \emph{l'impôt}
  \item \emph{Say it:} \emph{l'impôt}, once more
  \item \emph{Say it:} say \emph{l'impôt}
  \item \emph{Recall:} say the French for the truth, then the French for a lie, then say \emph{l'impôt} again
\end{itemize}

\begin{tcolorbox}[breakable,colback=teal!4,colframe=teal!35!black,title={Before you move on}]
What does l'impôt mean? (\textquotedblleft{}A tax\textquotedblright{}.) Say it once more.
\end{tcolorbox}
\section[le formulaire]{le formulaire — a form}

\label{lesson:FR-C252-formulaire}



\begin{quote}
Before the new one: say the French for a lie, then the French for a tax.
\end{quote}

\subsection*{You'll want to know: le formulaire}
\textbf{le formulaire} — \textquotedblleft{}a form\textquotedblright{}.

\begin{cousinweb}
\textbf{formulaire} $\leftarrow$ Latin \emph{formularium}, \textquotedblleft{}a book of forms\textquotedblright{}, from \emph{formula}, \textquotedblleft{}a little form\textquotedblright{} — a learned borrowing. (EtymDB 2.1)
\end{cousinweb}

\begin{grammarlens}[title={What you've built}]
One more word for this chapter.
\end{grammarlens}

\subsection*{Your turn}
\begin{itemize}
\raggedright
  \item \emph{Say it:} \emph{le formulaire}
  \item \emph{Say it:} \emph{le formulaire}, once more
  \item \emph{Say it:} say \emph{le formulaire}
  \item \emph{Recall:} say the French for a lie, then the French for a tax, then say \emph{le formulaire} again
\end{itemize}

\begin{tcolorbox}[breakable,colback=teal!4,colframe=teal!35!black,title={Before you move on}]
What does le formulaire mean? (\textquotedblleft{}A form\textquotedblright{}.) Say it once more.
\end{tcolorbox}
\section[l'enveloppe]{l'enveloppe — an envelope}

\label{lesson:FR-C252-enveloppe}



\begin{quote}
Before the new one: say the French for a tax, then the French for a form.
\end{quote}

\subsection*{You'll want to know: l'enveloppe}
\textbf{l'enveloppe} — \textquotedblleft{}an envelope\textquotedblright{}.

\begin{cousinweb}
\textbf{enveloppe} — formed from the verb for \textquotedblleft{}to wrap up\textquotedblright{} (EtymDB 2.1), built from \emph{en-} and an Old French verb for \textquotedblleft{}to wrap\textquotedblright{} of uncertain origin.
\end{cousinweb}

\begin{grammarlens}[title={What you've built}]
One more word for this chapter.
\end{grammarlens}

\subsection*{Your turn}
\begin{itemize}
\raggedright
  \item \emph{Say it:} \emph{l'enveloppe}
  \item \emph{Say it:} \emph{l'enveloppe}, once more
  \item \emph{Say it:} say \emph{l'enveloppe}
  \item \emph{Recall:} say the French for a tax, then the French for a form, then say \emph{l'enveloppe} again
\end{itemize}

\begin{tcolorbox}[breakable,colback=teal!4,colframe=teal!35!black,title={Before you move on}]
What does l'enveloppe mean? (\textquotedblleft{}An envelope\textquotedblright{}.) Say it once more.
\end{tcolorbox}
\section[le timbre]{le timbre — a stamp}

\label{lesson:FR-C252-timbre}



\begin{quote}
Before the new one: say the French for a form, then the French for an envelope.
\end{quote}

\subsection*{You'll want to know: le timbre}
\textbf{le timbre} — \textquotedblleft{}a stamp\textquotedblright{}.

\begin{cousinweb}
\textbf{timbre} $\leftarrow$ Old French \emph{timbre}, \textquotedblleft{}a drum, a bell\textquotedblright{} $\leftarrow$ Latin \emph{tympanum}, from Greek \emph{tympanon}, \textquotedblleft{}a drum\textquotedblright{} — inherited (EtymDB 2.1): then the mark struck on a letter.
\end{cousinweb}

\begin{grammarlens}[title={What you've built}]
One more word for this chapter.
\end{grammarlens}

\subsection*{Your turn}
\begin{itemize}
\raggedright
  \item \emph{Say it:} \emph{le timbre}
  \item \emph{Say it:} \emph{le timbre}, once more
  \item \emph{Say it:} say \emph{le timbre}
  \item \emph{Recall:} say the French for a form, then the French for an envelope, then say \emph{le timbre} again
\end{itemize}

\begin{tcolorbox}[breakable,colback=teal!4,colframe=teal!35!black,title={Before you move on}]
What does le timbre mean? (\textquotedblleft{}A stamp\textquotedblright{}.) Say it once more.
\end{tcolorbox}
\section[le colis]{le colis — a parcel}

\label{lesson:FR-C252-colis}



\begin{quote}
Before the new one: say the French for an envelope, then the French for a stamp.
\end{quote}

\subsection*{You'll want to know: le colis}
\textbf{le colis} — \textquotedblleft{}a parcel\textquotedblright{}.

\begin{cousinweb}
\textbf{colis} $\leftarrow$ Italian \emph{colli}, \textquotedblleft{}loads\textquotedblright{}, the plural of \emph{collo}, \textquotedblleft{}a neck, a load carried on it\textquotedblright{}, $\leftarrow$ Latin \emph{collum}, the Latin of \emph{cou} — borrowed.
\end{cousinweb}

\begin{grammarlens}[title={What you've built}]
One more word for this chapter.
\end{grammarlens}

\subsection*{Your turn}
\begin{itemize}
\raggedright
  \item \emph{Say it:} \emph{le colis}
  \item \emph{Say it:} \emph{le colis}, once more
  \item \emph{Say it:} say \emph{le colis}
  \item \emph{Recall:} say the French for an envelope, then the French for a stamp, then say \emph{le colis} again
\end{itemize}

\begin{tcolorbox}[breakable,colback=teal!4,colframe=teal!35!black,title={Before you move on}]
What does le colis mean? (\textquotedblleft{}A parcel\textquotedblright{}.) Say it once more.
\end{tcolorbox}
